\documentclass[11pt,a4paper]{article}
\usepackage[margin=22mm]{geometry}
\usepackage{amsmath,amssymb,mathtools}
\usepackage[hidelinks]{hyperref}
\usepackage{booktabs,array,microtype}
\setlength{\parindent}{1.5em}
\setlength{\parskip}{0.35em}
\newcommand{\floor}[1]{\left\lfloor #1\right\rfloor}
\newcommand{\ceil}[1]{\left\lceil #1\right\rceil}
\newcommand{\Z}{\mathbb Z}
\DeclareMathOperator{\lcm}{lcm}

\title{A Four-Point-Line Construction for OEIS A172992}
\author{}
\date{September 6, 2026}

\begin{document}
\maketitle
\vspace{-2.5em}

\begin{abstract}
We verify a construction that places $n$ rational points on four supporting
lines.  The correct parameter is $v=\ceil{(s-d-1)/2}$ (with $d$, not $c$).
With this correction, for every $n\ge28$ the construction guarantees
$\floor{(n^2+12)/28}$ lines containing exactly four selected points, and a
uniform dilation converts all coordinates to integers.  In particular, it
improves the known lower bounds for OEIS A172992 at $n=59,60$ to $124,129$.
These are lower bounds, not proofs of optimality.
\end{abstract}

\section*{Construction and collinearity}
Set
\[
 s=\floor{\frac{3(n+1)}7},\quad
 a=\floor{\frac s2},\quad b=\ceil{\frac s2},\quad
 c=\floor{\frac{n-s}2},\quad d=\ceil{\frac{n-s}2},
\]
and
\[
 u=\floor{\frac{s-c-1}{2}},\qquad
 v=\ceil{\frac{s-d-1}{2}}.
\]
Define the integer intervals
\[
\begin{split}
 A&=[0,a-1]\cap\Z,\qquad B=[0,b-1]\cap\Z,\\
 C&=[u,u+c-1]\cap\Z,\\
 D&=[-(b-1)+v,\;-(b-1)+v+d-1]\cap\Z.
\end{split}
\]
On $Y=0,X=0,Y=X,Y=-X$, respectively, place
\[
\begin{aligned}
 P_x&=\left(\frac1{3n+x},0\right),&
 Q_y&=\left(0,\frac1{n+y}\right),\\
 R_z&=\left(\frac1{4n+z},\frac1{4n+z}\right),&
 S_w&=\left(\frac1{2n+w},-\frac1{2n+w}\right).
\end{aligned}
\]
There are $a+b+c+d=s+(n-s)=n$ points.  The line through $P_x,Q_y$ is
\[
 (3n+x)X+(n+y)Y=1.
\]
Substitution of $R_z$ and $S_w$ gives, respectively,
\[
 \frac{4n+x+y}{4n+z}=1\iff z=x+y,
 \qquad
 \frac{2n+x-y}{2n+w}=1\iff w=x-y.
\]
Thus, whenever $x\in A,y\in B,x+y\in C,x-y\in D$, the four points
$P_x,Q_y,R_{x+y},S_{x-y}$ are collinear.  Such a line has a unique
intersection with each supporting line, so it contains exactly these four
selected points.  Distinct pairs $(x,y)$ give distinct lines because their
two axis intercepts are different.

\section*{Counting}
Put $L=-(b-1)+v$.  For fixed $x\in A$, the admissible values of $y$ form an
integer interval.  Hence the number of guaranteed four-point lines is
\[
 N(n)=\sum_{x=0}^{a-1}
 \Bigl[\min\{b-1,u+c-1-x,x-L\}
 -\max\{0,u-x,x-L-d+1\}+1\Bigr]_+,
 \tag{1}
\]
where $[t]_+=\max(t,0)$.  Formula (1) is also a direct, geometry-free way to
check the count.

Write $n=28q+r$, $0\le r<28$.  Substituting the endpoints in (1) and
simplifying separately in the 28 residue classes gives
\[
 N(n)=28q^2+2qr+h_r,\qquad h_r=\floor{\frac{r^2+12}{28}}.
\]
For auditability, all values of $h_r$ are displayed below; the calculation
only compares the finitely many integer-interval endpoints in (1).
\begingroup\scriptsize
\[
\begin{array}{c|rrrrrrrrrrrrrr}
r&0&1&2&3&4&5&6&7&8&9&10&11&12&13\\ \midrule
h_r&0&0&0&0&1&1&1&2&2&3&4&4&5&6\\[2pt]
r&14&15&16&17&18&19&20&21&22&23&24&25&26&27\\ \midrule
h_r&7&8&9&10&12&13&14&16&17&19&21&22&24&26
\end{array}
\]
\endgroup
Consequently,
\[
 N(n)=\floor{\frac{(28q+r)^2+12}{28}}
      =\floor{\frac{n^2+12}{28}},\qquad n\ge28.
\]

\section*{Integer coordinates and the cases 59 and 60}
All coordinates are rational.  Let $M$ be the least common multiple of all
positive denominators in
\[
 \{3n+x:x\in A\}\cup\{n+y:y\in B\}\cup
 \{4n+z:z\in C\}\cup\{2n+w:w\in D\}.
\]
The dilation $(X,Y)\mapsto(MX,MY)$ produces integer coordinates.  Being
invertible and linear, it preserves collinearity and the number of selected
points on each line.

For $n=59$ the parameters are
$(s,a,b,c,d,u,v)=(25,12,13,17,17,3,4)$, and (1) gives $124$.
For $n=60$ they are $(26,13,13,17,17,4,4)$, and (1) gives $129$.
As checked on September 6, 2026, the OEIS A172992 page still lists
$a(59)=123,a(60)=126$ and notes that its terms are the best solutions known,
not necessarily optimal ones.\footnote{OEIS A172992,
\url{https://oeis.org/A172992}, accessed September 6, 2026.}
The rigorous conclusion is therefore
\[
 A172992(59)\ge124,\qquad A172992(60)\ge129.
\]
If the sequence is updated according to its convention of recording the
best known constructions, the candidate entries are $124$ and $129$; this
construction does not rule out still larger values.

\end{document}
